\scanpage{24}
satisfaisant une formule de transitivité, % L1
et formule de projection % L2
% slanted
\margin{Considérer aussi la formule de dualité entre $\mathbb R f_*,f^*$, et les foncteurs $f^!,\mathbb R f_!$ et leurs relations d'adjonction et enfin les formules de \unsure{Künneth|chgt de base|standard} reliant tous ces foncteurs .....} % L2a
\[
\mathbb R f_*(M\otimes\mathbb L f^*(N))\simeq\mathbb R \unsure{f|f_*}(M)\otimes N % L3
\]

3) Supposons $X=\varprojlim X_i$, système projectif % L4
filtrant \unsure{en. affine|ess. affine}. Alors par les foncteurs % L5
images inverses, on a % L6
\[
\mathcal M^+(X)\simeq\underset{i}{\underset{\longrightarrow}{\mathrm{Lim}}}\,\mathcal M^+(X_i) % L7
\]

En particulier, si $X$ est de type fini % L8
sur $S=\mathrm{Spec}(A)$, alors $X$ est limite de % L9
préschémas $X_i$ de type fini sur $\mathbb Z$, et % L10
la détermination de $\mathcal M^+(X)$ avec ses % L11
structures déjà envisagées est \deleted{\unsure{dét|dit}} \unsure{ramené|ramenée} % L12
au cas des préschémas de type fini. % L13

De $=$, si $(S_i)$ est système % L14
projectif filtrant \unsure{en. affine|ess. affine}, $S=\varprojlim S_i$. % L15
et si $X,Y$ de t. f. sur $S$ sont % L16
définis par $(X_i)$, $(Y_i)$ de la % L17
façon habituelle, si on part de % L18

% archive sheet number: 23
